\documentclass[literature.tex]{subfiles}

\begin{document}
\prose{Hobit, aneb Cesta tam a zase zpátky}{J. R. R. Tolkien}
\teorieA{
Román považovaný za základní kámen fantasy.
Psán pro dětského posluchače a subžánrově náleží do \texthl{high fantasy}.
}{
Důraz na konsistentní svět, ovlivněn \texthl{severskou a staroanglickou mytologií a folkrórem}.
}{
Kniha je napsána v angličtině, ale zahrnuje prvky vytvořených jazyků,
jako je elfština (Quenya, Sindarin, Telerin, Noldorin, Nandorin, Doriath, Ilkorin, Avarin)
a trpasličí řeč (Khuzdûl, Iglishmêk, QuasiKhuzdul), odrážející Tolkienovu lásku k lingvistice.
Slovní zásoba čerpá ze staroangličtiny a dalších historických jazyků, což přispívá k bohatosti textu,
přizpůsobené dětskému čtenáři.
}{}{}{
Přímý a věcný jazyk, ideální pro povídku s rychlým spádem, co má zaujmout mladší diváky.
Vyskytují se písně a humor, které dále odlehčují příběh, který by jinak mohl být temný.
}{
Nakonec to už pan Pytlík nevydržel,
\enquote{Hrom aby do tebe, Šmaku, ty nestvůro!} vypískl hlasitě.
\enquote{Přestaň si hrát na schovávanou!
Posviť mi, a pak si mě třeba sežer, jestli mě dokážeš chytit!}\par
Síní, kterou nebylo vidět, se rozlehla slabá ozvěna, ale neozvala se žádná odpověď.\par
Bilbo vstal a shledal, že neví, kterým směrem se má obrátit.
\enquote{Teď bych rád věděl, na co si ten Smak u všech všudy hraje,} řekl si. 
\enquote{Určitě dnešního dne (nebo dnešní noci, či co vlastně je) není doma.
Jestli Oin a Gloin neztratili svá křesadla a troud,
snad bychom si mohli trochu posvítit a rozhlédnout se kolem, než se zas obrátí štěstí.}\par
\enquote{Světlo!} vykřikl. \enquote{Může někdo udělat světlo?}\par
Trpaslíci se samozřejmě hrozně vyplašili, když Bilbo spadl s žuchnutím do síně,
a zůstali schouleně sedět před koncem tunelu, kde je zanechal.\par
\enquote{Pst! Pst!} syčeli, když uslyšeli jeho hlas,
a třebaže to hobitovi pomohlo určit, kde jsou, trvalo nějaký čas,
než z nich dostal něco jiného.
Ale nakonec, když začal skutečně dupat na zem a ječel \enquote{světlo!},
co mu jen stačil pronikavý hlas, Thorin se podvolil
a poslal Oina a Gloina zpátky k jejich rancům na horním začátku tunelu.\par
}
\teorieB{
\wtem{Bilbo Pytlík}{Hobit~z Kraje milující pohodlí,
zdráhavě se přidává~k výpravě trpaslíků.
Nachází kouzelný prsten~a odhaluje svou odvahu.
Z bázlivého se stává hrdinou, ctí domov~i dobrodružství.}
\wtem{Gandalf}{Moudrý čaroděj, Šedý poutník, organizuje výpravu~k Osamělé hoře.
Vede skupinu, zasahuje v~klíčových chvílích, zůstává tajemný.
Vidí Bilbův potenciál~a je morálním kompasem.}
\wtem{Thorin Pavéza}{Vůdce trpaslíků, dědic království, touží po Šmakově pokladu.
Hrdý~a loajální, ale chamtivost ho zraňuje.
Umírá~v Bitvě pěti armád po usmíření~s Bilbem.}
\wtem{Šmak}{Drak střežící poklad Osamělé hory, arogantní~a chytrý.
Jeho pýcha odhalí slabinu, útočí na Jezerní město.
Zabije ho Bard, což je dějový vrchol.}
\wtem{Glum}{Zkažený tvor, dříve hobit, posedlý prstenem, žije pod Mlžnými horami.
Setkání~s Bilbem~a hádanky jsou klíčové.
Ztělesňuje varování před mocí.}
\wtem{Elrond}{Moudrý elf, pán Roklinky, pomáhá trpaslíkům~s mapou.
Vyzařuje klid~a učenost, symbolizuje stabilitu.
Roklinka je bezpečným útočištěm.}
\wtem{Bard}{Lučištník z~Jezerního města, zabije Šmaka~a stává se vůdcem.
Skromný, ale odvážný, představuje lidskou naději~a obnovu.}
\wtem{Beorn}{Měňavec, medvěd-člověk, žije osaměle.
Pomáhá trpaslíkům útočištěm~a~v Bitvě pěti armád bojuje proti goblinům.}
\wtem{Dwalin}{Loajální~a bojovný trpaslík, věrný Thorinovi, jeden~z prvních~u Bilba.}
\wtem{Balin}{Moudrý trpaslík, přítel Bilba, známý~z „Pána prstenů“.}
\wtem{Fili~a Kili}{Thorinovi synovci, mladí, hraví~a stateční,
sdílejí dobrodružný duch.}
\wtem{Oin~a~Gloin}{Bratři, Oin zná léky, Gloin je houževnatý,
otec Gimliho~z „Pána prstenů“.}
\wtem{Dori~a~Nori}{Dori je silný a~pomáhá Bilbovi, Nori tichý, ale šikovný,
oba věrní skupině.}
\wtem{Bofur, Bombur~a~Bifur}{Bofur je veselý, Bombur těžký~a komický,
Bifur tichý, všichni spolehliví.}
\wtem{Ori}{Mladý písař, zapisuje události, později známý~z Morie.}
}{
Bilbo Pytlík je starým čarodějem Gandalfem vhozen do pohádkového děje.
Postupně se učí začlenit a zlepšuje své schopnosti.
Nejde o změnu, stále je to Bilbo, ale naučí se žít i se svojí další stránkou.
Samotný podtitul to popisuje asi nejlépe, \enquote{Cesta tam a zase zpátky,} nejedná se o změnu životního stylu, ale o dobrodružnou vycházku.
Výjimečné vyčlenění z normálu.
}{
Struktura děje je lineární, rozdělena do 19 kapitol, každá posouvající příběh.
Odehrává se ve Středozemi.
Čas je situován do Třetího věku, přibližně v letech 2941–2942.
}{}
\historie{
Hobit byl napsán v 30. letech 20. století,
během meziválečného období, kdy v Evropě rostly totalitní režimy, zejména nacismus v Německu.
Tolkien vyjádřil odpor k těmto ideologiím, což mohlo ovlivnit témata boje dobra se zlem.
}{}{}{
Anglie po Druhé světové válce.
}{
Tolkien byl ovlivněn staroanglickými a severskými ságami,
jako je \texthl{Beowulf} a \texthl{Poetická Edda},
svými zkušenostmi z 1. světové války,
křesťanskou vírou a venkovským prostředím Anglie.
}{
Narodil se v Jižní Africe v roce 1892,
přestěhoval se do Anglie jako dítě, studoval na Oxfordu,
sloužil v 1. světové válce a stal se profesorem anglického jazyka a literatury na Oxfordu.
Napsal Hobita v 30. letech a Pána prstenů v 50. letech, zemřel v roce 1973
}{
Mezi další významná díla patří „Pán prstenů“ (třídílný román), „Silmarillion“ (sbírka mýtů),
„Nedokončené příběhy“, „Historie Středozemě“ (série knih) a různé básně a eseje
}
\kritika{
obit ve svém čase nebyl tak slavný, proslavila se až následující trilogie.\\
\vspace*{1pt}\\
Při prvním vydání v roce 1937 byl „Hobit“ přijímán jako dětský příběh,
s pozitivními recenzemi,
jako ocenění od New York Herald Tribune za nejlepší dětskou fikci roku 1938.
Postupem času se jeho status změnil, stal se klasikou fantasy literatury,
s uznáním jako „Nejdůležitější román 20. století pro starší čtenáře“
v anketě Books for Keeps v roce 2012.
}{
Na Hobita Tolkien navázal se svojí trilogií Pán prstenů, která dobila svět
a~stala se základním kamenem žánru fantasy.
Kdy prakticky každé fantasy má nějaké prvky společné s Tolkienovou tvorbou.\\
\vspace*{1pt}\\
„Hobit“ pomohl popularizovat žánr high fantasy,
stanovil konvence jako detailní budování světa, epické výpravy a mytologické prvky,
ovlivnil spisovatele jako Terry Pratchett a George R. R. Martin,
a přispěl k rozvoji moderní fantasy literatury.
}{
Témata jako důležitost domova a komunity, lákavost dobrodružství,
boj dobra se zlem a hodnota přátelství a loajality zůstávají relevantní.
V globalizovaném světě je hodnota komunity stále důležitá, touha po dobrodružství přetrvává,
a morální konflikty jsou nadčasové.
}
\nazor{
Pohádka co dala zrodit celý moderní žánr fantasy. Částečně historie, částečně základ archetypu, který je obsažen všude.
}
\explain{
\textbl{High fantasy --} subžánr fantasy, kde se děj odehrává v úplně vytvořeném světě, jako je Středozem, s epickými konflikty
}
\wpage
\end{document}